La administración de operaciones es la función encargada de gestionar los procesos que transforman recursos en bienes o servicios. Dentro de cualquier empresa, ya sea de manufactura o de servicios, esta función determina cómo se produce, con qué eficiencia y bajo qué estándares de calidad. Representa una de las áreas más importantes porque conecta directamente los objetivos estratégicos de la organización con la ejecución diaria de sus actividades \parencite{heizer2014}.

En términos prácticos, la administración de operaciones no es exclusiva de las fábricas. Un hospital que gestiona camas, personal y turnos; una empresa de logística que coordina rutas y entregas; o un restaurante que programa ingredientes y cocineros, todos aplican principios de administración de operaciones. Esto evidencia que su alcance es amplio y su relevancia, independiente del tamaño o sector de la organización.

Para los gerentes y responsables de una empresa, entender la administración de operaciones significa comprender cómo se crean los productos o servicios que la empresa ofrece. Cuando se conoce ese proceso, es más fácil reducir costos, mejorar la calidad, atender mejor al cliente y adaptarse a los cambios del mercado. Por eso, esta disciplina forma parte esencial de la formación gerencial moderna \parencite{chase2014}.

El propósito de esta investigación es explorar la administración de operaciones, identificando sus conceptos clave, su relevancia organizacional y las principales decisiones que los gerentes deben tomar para gestionar operaciones de manera eficiente.